\documentclass[10pt]{article}

\usepackage[utf8]{inputenc}

\usepackage[T1]{fontenc}

\usepackage{graphicx}

\usepackage[export]{adjustbox}

\graphicspath{ {./images/} }

\usepackage{amsmath}

\usepackage{amsfonts}

\usepackage{amssymb}

\usepackage[version=4]{mhchem}

\usepackage{stmaryrd}



\begin{document}

но основной вывод, содержащийся в формуле (4), сохраняет свою силу.



Существует целый ряд примеров получения плодотворных выводов и установления целесообразных удобных методов обработки әкспериментальных результатов с использованием Р. а. и л-теоремы. Особая польза достигается за счет сокращения числа аргументов у искомых функций; это сокращение обеспечи-



\includegraphics[max width=\textwidth, center]{2023_07_08_df28f13285e28442c3efg-1}

основании постановки задачи с использованием опытных данных или нек-рых гипотез о механизмах существенных эффектов в изучаемых вопросах.



Основная практич. польза состоит в установлении возможности перенесения результатов опыта в одних условиях на др. условия, в к-рых опыт не проводился. Напр., по опытам с водой, движущейся в трубе с фиксированным диаметром, можно давать автоматически нужные ответы о движении воды в трубах той же формы, но других размеров, о движении в трубах нефти, воздуха и т. п. Однако для возможности перенесения результатов данного опыта на др. опыты необходимо обеспечивать в соответствующих случаях одинаковые значения безразмерных параметров. Значения нек-рых из этих величив легко обеспечиваются геометрич. условиями постановки опытов. Значения других параметров связаны с физико-механич. условиями проведения опытов.



Примерами применения Р. а. являются примеры установления из постановки задач автомодельности искомых решений. Свойство а в т о м о д е л в в о с т и состоит в следующем. При первом подходе к математич. разрешению ставящейся задачи можно отметить независимые переменные аргументы функций; обычно это три координаты точек в нек-ром объемө пространства и время. Кроме этих переменных величин среди аргументов искомых функций могут фигурировать век-рые задаваемые постояныые параметры, возникающие при описании свойств изучаемой модели и при выделении конкретной схематизированной задачи. Всякую задачу можно сформулировать в безразмерных переменных как для искомых функций, так и для их везависимо изменяющихся аргументов. Автомодельность проявляется в том, что сокращается число независимых безразмерных переменных аргументов по сравнению с числом размерных.



Эффективные результаты, связанные с автомодельностыо, получаются в теории нек-рых одномерных неустановившихся задач со сқерическими, цилиндрическими и плоскими волнами, когда имеются только две независимые переменные: координата $r$ с размерностью длины и время $t$. Если из определяющих размерных постоянных, присутствуюцих в постановке задачи, нельзя образовать две комбинации - одну с размерностью длины, а другую с размерностью времени, то единственным иеременным безразмерным аргументом в искомых безразмерных функциях будет парамезр вида $\lambda=C^{\prime} r / t^{\alpha}$, где $\alpha$ - нек-рый показатель, а $C-$. размерная постоянная, выражающаяся терез заданные постоянные, фигурирующие в постановке задачи. В частности, если математич. задача сводилась к решению нек-рой системы дифференциальных уравнений с частными производными по $r$ и $t$, то при наличии только одного переменного параметра $\lambda$ уравнения с частными производными по $r$ и $t$ преобразуются к обыкновенным дифференциальным уравнениям с одной независимой переменной $\lambda$, вследствие чего математич. задача сильно упрощается. Таким путем было получено решение многих практически важных задач, напр. задачи о возмущенном движении воздуха, вызванном концентрированным атомным взрывом в атмосфере.



Соображения $\mathrm{P}$. а. для автомодельных процессев могут служить не только для сведения уравнений с qастными производными к обыкновенным, но и для получения конечных соотношений между искомыми функциями (интегралов этих уравнений), что позволяет находить решения нек-рых автомодельных задач в замкнутом формульном виде. Эти методы были продемонстрированы при рөшении задач о сильном взрыве в атмосфере и в задачах о распространении взрывных волн в газовых средах, в к-рых в исходном состоянии плотность переменна и изменяется вдоль радиуса по стөпенному закону. Эти задачи ставятся как сильно схематизированные модели взрыва звезд.



P. а. непосредственно связан с понятием о физич. подобии, к-рое изучается в подобии теории.



Лит.: [1] Б р и Д ж м е н П. В., Анализ размерностей, пер. с англ., Л.- - М., 1934; [2] С е д о в JI. И., Методы подобия и



РАЗМЕРНОСТИ АДДИТИВНЫЕ СВОЙСТВА свойства, выражающие связь размерности топологич. пространства $X$, представленног о в виде суммы своих подпространств $X_{\alpha}$, с размерностями пространств $X_{\alpha}$. Имеется несколько видов Р. а. с.



Т е о рем а с уммы. Если хаусдорфово и нормальное пространство $X$ представимо в виде конечной или счетной суммы своих замкнутых подмножеств $X_{i}$, то



\[

\operatorname{dim} X_{i}=\sup _{i} \operatorname{dim} X_{i}

\]



Если, дополнительно, пространство $X$ совершенно нормально или наследственно паракомпактно, то



\[

\text { Ind } X=\sup _{i} \text { Ind } X_{i} \text {. }

\]



Локально конечная те орема су мм ы. Если хаусдорфово и нормальное пространство $X$ представлено в виде суммы локально конечной системы своих замкнутых подмножеств $X_{\alpha}$, то



\[

\operatorname{dim} X=\sup _{\alpha} \operatorname{dim} X_{\alpha}

\]



Если, дополнительно, пространство $X$ соверпенно нормально или наследственно паракомпактно, то



\[

\text { Ind } X=\sup _{\alpha} \text { Ind } X_{\alpha} \text {. }

\]



Те о рем а сло ж н и я. Если пространство $X$ хаусдорфово, наслөдственно нормально и $X=A \cup B$, то



\[

\operatorname{dim} X \leqslant \operatorname{dim} A+\operatorname{dim} B+1

\]



(ф) о м ула Менге р а кроме того, пространство $X$ совершенно нормально, то



\[

\text { Ind } X \leqslant \text { Ind } A+\text { Ind } B+1 .

\]



Метрич пространство $R$ имеет размерность $\operatorname{dim} R \leqslant n$ тогда и только тогда, когда



$R=\bigcup_{i=1}^{n+1} R_{i}, \operatorname{dim} R_{i} \leqslant 0, i=1, \ldots, n+1 ; n=0,1,2, \ldots$



В хаусдорфовом наследственно нормальном пространстве $X$ для любого замкнутого подмножества $F$ вылолняются равенства



\[

\begin{aligned}

& \operatorname{dim} X=\max (\operatorname{dim} F, \operatorname{dim} X \backslash F), \\

& \text { Ind } X=\max (\text { Ind } F, \text { Ind } X \backslash F) .

\end{aligned}

\]



Б. А. Пасынков. - Ро тасть топологии, в к-рой для каждого компакта, а впоследствии и для более общих классов топологич. пространств тем или иным өстөственным образом определяется числовой топологич. инвариант - размерность, совпадающий, если $X$ есть полиэдр (в частности, многообразие), с его числом измерений в смысле әлементарной или дифференциальной геометрии. Первое общее определение раз-





\end{document}